% Paper 6 measured-evidence section: same-stack variant deltas from N2 tensors
\section{Measured Evidence: N2 Same-Stack Variant Deltas}
\label{sec:p6-measured}

The schema and entries answer \emph{what the stack is}; the question for an
empirical paper is \emph{what the variant label actually buys, once the
stack is held fixed}. We re-executed the four GRPO-family methods
(\texttt{grpo}, \texttt{aero}, \texttt{gift}, \texttt{areal}) on a single
identical stack---Tinker-managed sampler, $G{=}8$, seed $0$, $40$ steps,
$16$ prompts per step---logging the full per-prompt reward tensor at every
step (\texttt{experiments/results/n2\_reward\_tensor\_resume/}).

\subsection{Schema validation}
At registry commit \texttt{ca5b792}, all 46 shipped records (28 \texttt{stack},
18 \texttt{variant\_delta}) validate against \texttt{schema.json}; the two
amendment sidecars do not, because they carry no \texttt{record\_type} or
\texttt{schema\_version} (Appendix~\ref{sec:p6-appendix-layout}). An earlier
run of the CI-style auditor (\texttt{scripts/p5p8/registry\_validate.py},
\texttt{experiments/results/p5p8/registry\_schema\_check.tsv}) passed
31/31 records (20 stack, 11 variant-delta) on a smaller snapshot.

\subsection{Measured deltas on the same stack}
Table~\ref{tab:p6-measured-deltas} reports paired bootstrap means and
$95\%$ CIs on the last-$10$ steps (deterministic LCG, $n_{\text{boot}}=2000$).
The per-prompt pooled view pools each prompt's mean reward over the steps
it sees. It describes prompt variation within one seed, not a seed-level
counterfactual or a causal algorithm effect.

\begin{table}[t]
\centering
\caption{Same-stack variant deltas ($\Delta = \text{variant} - \texttt{grpo}$)
from the N2 four-method run (seed $0$, $G{=}8$, last-10 steps; paired
bootstrap CI on per-step scalar, then per-prompt pooled). The reward
deltas are all $\leq 2$pp in magnitude; for the AERO and AREAL reward
\emph{decreases} the step-bootstrap CI excludes $0$, but that interval
resamples 10 autocorrelated steps of one seed and is not a seed-level
significance test. The largest effects
are on the length distribution: AERO and AREAL \emph{lengthen}
completions by $\sim$31 tokens and reduce length-CV. GIFT raises ZVF
($+0.125$) and shifts the loss by its gamma-style additive baseline
constant.}
\label{tab:p6-measured-deltas}
\small
\resizebox{\linewidth}{!}{%
\begin{tabular}{lllrrr}
\toprule
Metric (scope)        & variant & grpo mean & var. mean & $\Delta$ & 95\% CI \\
\midrule
reward\_mean (last-10) & aero    & 0.861 & 0.847 & $-0.014$  & $[-0.023,\,-0.005]^{\dagger}$ \\
reward\_mean (last-10) & gift    & 0.861 & 0.877 & $+0.016$  & $[-0.007,\,+0.040]$ \\
reward\_mean (last-10) & areal   & 0.861 & 0.841 & $-0.020$  & $[-0.032,\,-0.008]^{\dagger}$ \\
zvf (last-10)          & gift    & 0.775 & 0.900 & $+0.125$  & $[+0.081,\,+0.175]^{\dagger}$ \\
mean\_len (last-10)    & aero    & 319.6 & 350.7 & $+31.08$  & $[+27.38,\,+35.09]^{\dagger}$ \\
mean\_len (last-10)    & areal   & 319.6 & 349.7 & $+30.13$  & $[+25.17,\,+34.60]^{\dagger}$ \\
cv\_len (last-10)      & aero    & 0.325 & 0.291 & $-0.034$  & $[-0.069,\,-0.014]^{\dagger}$ \\
cv\_len (last-10)      & areal   & 0.325 & 0.287 & $-0.038$  & $[-0.062,\,-0.019]^{\dagger}$ \\
loss (last-10)         & gift    & $+2.3\!\times\!10^{2}$ & $-1.6\!\times\!10^{4}$ & $-1.67\!\times\!10^{4}$ & $[-1.94\!\times\!10^{4},\,-1.49\!\times\!10^{4}]^{\dagger}$ \\
\midrule
reward\_mean (per-prompt, $n{=}16$) & aero  & --- & --- & $-0.007$ & $[-0.050,\,+0.028]$ \\
reward\_mean (per-prompt, $n{=}16$) & gift  & --- & --- & $+0.011$ & $[-0.013,\,+0.031]$ \\
reward\_mean (per-prompt, $n{=}16$) & areal & --- & --- & $-0.006$ & $[-0.044,\,+0.016]$ \\
\bottomrule
\end{tabular}%
}\\
\footnotesize{$^{\dagger}$ Step-bootstrap CI excludes $0$ (single seed);
per-prompt rows report only the pooled contrast over $n{=}16$ prompts. Source: \texttt{experiments/results/n2\_reward\_tensor\_resume/}.
}
\end{table}

\subsection{What this rules in and out}
These intervals summarize one logged trajectory per method. Adjacent steps
are not independent training replicates; significance of a step bootstrap
does not establish a reproducible method effect. No cross-seed reward
ranking or mechanistic interpretation follows. The last-10 window is part
of the estimand, not a full-trajectory result. The run is also near
ceiling: over the last 10 steps, $72.5\%$ of GRPO's groups are all-correct
(\texttt{frac\_all\_one}; $82.5\%$ for GIFT), so most of the measured ZVF
reflects mastered prompts, and the run cannot by itself support a claim that
the variant label does not move outcomes once the stack is fixed.
\endinput % Historical stronger interpretation preserved below.
The measured reward deltas are all $\leq 2$pp in magnitude at the same
stack---direct quantitative corroboration of the Pillar-1 ``algorithm
barely moves outcome once stack is fixed'' reading. Two of the three
(AERO $-0.014$ and AREAL $-0.020$) are statistically significant on the
last-10 per-step scalar (CI excludes $0$; corroborated by the iter-190
raw-recompute audit in Sec.~\ref{sec:p6-iter190-raw-recompute}), but both are tiny
reward \emph{decreases}, not improvements, so the magnitude thesis holds
while the ``all within noise'' claim does not. The length / length-CV
deltas for AERO and AREAL are the largest same-stack effects, but they
are \emph{not} reward improvements: they are different
prompt--completion distributions under the same reward, with AERO and
AREAL \emph{lengthening} completions ($\sim$31 tokens) and tightening
the length-CV, consistent with sequence-level importance-ratio methods
reshaping the length distribution rather than the outcome. GIFT's loss
additive shift is a baseline constant, not an outcome.

The registry's job is to make these distinctions visible. An entry that
claims a ``DAPO'' outcome without a per-component delta status makes it
impossible to know whether the reported number reflects an algorithm
change or a stack change; the N2 row gives the registry a measured
reference for the same-stack baseline against which future entries can
be compared.
